% ---- pg-311 （书页 299） ----

\subsection{线性回归中的其他问题}

我们来简单讨论一下回归中的几个重要问题. 这些问题在回归分析的专著中得到了很大篇幅的讨论(见文献[10], [11], [13]和[23]).

\noindent\textbf{从一个大集合中选择预测变量}\quad 实践中往往很难立刻就得到一个合适的回归函数. 哪些预测变量应该被包含在模型中? 回归函数应取何种形式?

当可能的预测变量非常多时, 并非所有变量都应进入模型. 现在很容易找到选择``最佳''预测量子集的方法及有关计算机程序. 一类好办法是试用所有子集: 只用 $z_1$, 只用 $z_2$, 用 $z_1$ 和 $z_2$, $\cdots$, 等等. 最佳选择通过检查某个准则来确定, 常用的一个选择量是 $R^2$[见\hyperref[eq:7-9]{式(7-9)}]. 但是, $R^2$ 总是随着模型中变量的增加而增加. 虽然这个问题可以用修正的 $R^2$ 即 $\bar{R}^2=1-(1-R^2)(n-1)/(n-r-1)$ 来缓解, 但更好的变量选择准则可能应该是马罗(Mallow)的 $C_p$ 统计量(见文献[12]):
\[
C_p=\left(\frac{\left(\begin{array}{c}\text{包括常数项在内含 }p\text{ 个参数}\\[-0.3ex]\text{的子模型的残差平方和}\end{array}\right)}{(\text{全模型的残差方差})}\right)-(n-2p)
\]
对每个预测量子集作 $(p,C_p)$ 标绘图, 此图将显示出能很好预测响应观测值的模型. 典型的情况是: 好模型的 $(p,C_p)$ 坐标在 $45^\circ$ 线附近. 在\hyperref[fig:7.4]{图 7.4} 中, 我们用小圆圈来表示``最佳''预测变量子集.

\begin{mvfigure}
\begin{tikzpicture}[x=0.62cm,y=0.62cm,font=\small,line width=0.5pt]
% 坐标轴（带断轴）
\draw[->] (0,0)--(0,11.05);
\draw[->] (0,11.75)--(0,14.5);
% 断轴符号
\draw (-0.12,11.28)--(0.12,11.20)--(-0.12,11.45)--(0.12,11.37);
% 轴上的横虚线
\draw[dashed] (0.1,11.35)--(7.4,11.35);
% 45 度线
\draw (0,0)--(6,6);
% 横轴
\draw[->] (0,0)--(7.6,0) node[below right,inner sep=1pt] {$p=r+1$};
% 横轴刻度
\foreach \i in {1,...,7}{\draw (\i,0)--(\i,-0.18) node[below,inner sep=1pt] {\i};}
% 纵轴刻度（下半段）
\foreach \i in {1,...,11}{\draw (0,\i)--(-0.18,\i) node[left,inner sep=2pt] {\i};}
% 纵轴刻度（上半段）
\foreach \v/\lab in {11.75/{$1\,200$},12.8/{$1\,600$},13.85/{$1\,800$}}{\draw (0,\v)--(-0.18,\v) node[left,inner sep=2pt] {\lab};}
\node[left,inner sep=4pt] at (-0.6,7.0) {$C_p$};
% 数据点
\fill (1.15,13.85) circle (2.3pt); \node[right,inner sep=2pt] at (1.15,13.85) {$(0)$};
\fill (2.05,12.85) circle (2.3pt); \node[right,inner sep=2pt] at (2.05,12.85) {$(3)$};
\fill (1.45,12.35) circle (2.3pt); \node[right,inner sep=2pt] at (1.45,12.35) {$(2)$};
\fill (2.95,12.15) circle (2.3pt); \node[right,inner sep=2pt] at (2.95,12.15) {$(2,3)$};
\fill (2.90,10.55) circle (2.3pt); \node[right,inner sep=2pt] at (2.90,10.55) {$(1,3)$};
\fill (1.90,8.50) circle (2.3pt); \node[right,inner sep=2pt] at (1.90,8.50) {$(1)$};
\fill (3.05,2.60) circle (2.3pt); \draw (3.05,2.60) circle (4.6pt);
\node[right,inner sep=3pt] at (3.05,2.60) {$(1,2)$};
\fill (3.70,4.00) circle (2.3pt); \node[right,inner sep=2pt] at (3.70,4.00) {$(1,2,3)$};
\node[align=center] at (4.6,1.7) {括号中的数字代表\\预测变量的编号};
\end{tikzpicture}
\phantomsection\label{fig:7.4}\mvcap{图 7.4\quad 例 7.6 中计算机数据的 $\mathbf{C}_p$ 标绘图}
\mvcap{($z_1=$命令数, $z_2=$增删项目数, $z_3=$项目数)}
\end{mvfigure}

如果预测变量的数目非常多, 费用方面的考虑使能考察的模型个数受到限制. 另一种方法称为逐步回归法(见文献[13]), 这种方法力图选择重要的预测变量而不考虑所有可能性. 这种方法可用下列基本步骤来描述:

\noindent\textbf{步骤 1.}\quad 考虑所有可能的简单线性回归. 能够解释大部分 $Y$ 变差的预测变量(即与响应 $Y$ 相关最强的变量)是应首先被列入回归函数的变量.

\noindent\textbf{步骤 2.}\quad 下一个应被列入模型的变量是模型外变量中对回归平方和作出最显著贡献的那个变量. 贡献的显著性由 $F$ 检验确定(见结论 7.6). 在一个变量的贡献被认为显著之前必须超过的那个 $F$ 统计量的值, 常被称为 \textbf{$F$ 临界值}.

\noindent\textbf{步骤 3.}\quad 一旦又有某个变量进入了模型, 则应利用 $F$ 检验来检验回归方程中每一个原有变量对回归平方和贡献的显著性. 若 $F$ 统计量小于某个与指定显著性水平相对应的临界值, 就应将相应变量从回归函数中剔除.

\noindent\textbf{步骤 4.}\quad 将步骤 2 与步骤 3 不断重复, 直到所有可能增添的变量均不显著, 所有可能删除的变量均显著时为止.

由于整个过程是一步一步进行的, 不能保证这种方法能选出比如说三个``最佳''预测变
